% !TeX encoding = UTF-8
\documentclass[variant=guiaresuelta,cover=institutional,mode=screen]{ua-engineering-v6-article}
% !TeX encoding = UTF-8
\UASetUniversity{Universidad Austral}
\UASetFaculty{Facultad de Ingeniería}
\UASetDepartment{Departamento de Computación}
\UASetCampus{Pilar}
\UASetCareer{Ingeniería Informática}
\UASetCommission{Comisión A}
\UASetProfessor{Equipo docente}
\UASetSemester{Segundo cuatrimestre 2026}
\UASetCourse{Sistemas Distribuidos}
\UASetDate{2026}

\UASetSemester{Segundo cuatrimestre 2026}
\UASetDate{2026}

\UASetClassNumber{06}
\UASetClassTitle{Storage distribuido y almacenamiento permanente}
\UASetDocKind{Guía resuelta}
\title{Clase 06 --- Guía resuelta: persistencia, concurrencia y rendimiento bajo fallas}
\author{\UAProfessor}
\date{\UADate}

\begin{document}
\setlength{\emergencystretch}{2em}
\maketitle

\section{Cómo leer las soluciones}

Las respuestas son modelos de razonamiento, no recetas únicas. Una alternativa es válida si declara el modelo de falla, preserva el invariante, hace explícita la frontera de garantía y propone evidencia capaz de contradecir la elección.

\begin{uakeybox}
La secuencia esperada es: contexto $\rightarrow$ garantía $\rightarrow$ mecanismo $\rightarrow$ ejecución adversarial $\rightarrow$ costo $\rightarrow$ evidencia.
\end{uakeybox}

\section{Resoluciones}

\begin{uaexercise}
\textbf{1. Del problema general a los mecanismos.}
\end{uaexercise}
\begin{uaanswer}
Una tabla posible es:

\begin{center}
\small
\begin{tabular}{p{0.19\textwidth}p{0.25\textwidth}p{0.25\textwidth}p{0.20\textwidth}}
\toprule
Responsabilidad & Escenario sin cubrir & Garantía buscada & Mecanismo candidato \\
\midrule
Registrar decisiones & Compra confirmada desaparece tras crash & Todas y solo las confirmadas sobreviven & WAL y recovery \\
Controlar concurrencia & Reporte mezcla estados o bloquea checkout & Vista coherente con escritores activos & MVCC y snapshot \\
Elegir ruta física & Se recorren un millón de filas para devolver 20 & Trabajo proporcional al subconjunto útil & Índice y planificador \\
Optimizar escrituras & Muchas actualizaciones aleatorias saturan E/S & Ingestión sostenible & Árbol B o LSM según carga \\
Proteger capacidad & Compactación queda atrás y el backlog crece & Cola interna acotada y p99 controlado & Compactación y write stall \\
\bottomrule
\end{tabular}
\end{center}

La idea central es que ningún mecanismo constituye por sí solo ``el storage''. Cada uno cubre una responsabilidad distinta y desplaza un costo diferente.
\end{uaanswer}

\begin{uaexercise}
\textbf{2. Siglas y términos sin ambigüedad.}
\end{uaexercise}
\begin{uaanswer}
Una respuesta lograda incluye, al menos:

\begin{description}
\item[WAL.] Write-Ahead Logging, registro anticipado de escritura: evidencia para recuperar un cambio se registra antes de depender de la página modificada.
\item[MVCC.] Multi-Version Concurrency Control, control de concurrencia multiversión: mantiene versiones y reglas de visibilidad por snapshot.
\item[LSM-tree.] Log-Structured Merge-tree, árbol de fusión estructurado como log: escribe primero en memoria y SSTables, y reorganiza mediante compactación.
\item[ACK.] Acknowledgement, acuse de recibo emitido cuando la garantía prometida ya se cumplió.
\item[LSN.] Log Sequence Number, posición ordenada dentro del WAL.
\item[SQL.] Structured Query Language, lenguaje de consulta y modificación de datos.
\item[DDL.] Data Definition Language, parte de SQL que define estructuras; \texttt{CREATE INDEX} es DDL.
\item[OLTP.] Online Transaction Processing, carga de transacciones cortas y concurrentes.
\item[SLO.] Objetivo cuantitativo de nivel de servicio.
\item[RPO.] Pérdida máxima de datos aceptable respecto de un punto de recuperación.
\item[RTO.] Tiempo máximo aceptable para recuperar el servicio.
\item[COMMIT.] Decisión definitiva del motor sobre los efectos de la transacción.
\item[Log durable.] Prefijo del WAL que sobreviviría a la falla declarada.
\item[Flush durable.] Volcado concluido que hizo persistentes los bytes requeridos.
\item[Snapshot.] Regla lógica que determina qué versiones son visibles.
\item[Serializabilidad.] Equivalencia con alguna ejecución secuencial de las transacciones.
\item[SSTable.] Archivo ordenado e inmutable de pares clave--valor.
\item[Merge ordenado.] Fusión lineal de runs ya ordenados para producir otro run ordenado.
\item[Filtro de Bloom.] Resumen probabilístico que responde ausencia definitiva o presencia posible.
\item[Write stall.] Backpressure que reduce o detiene escrituras para contener backlog.
\end{description}

La condición ``ninguna sigla sin definir'' evita frases circulares como ``LSM compacta SSTables con Bloom'' antes de explicar esos objetos.
\end{uaanswer}

\begin{uaexercise}
\textbf{3. Identidad lógica y objetos físicos.}
\end{uaexercise}
\begin{uaanswer}
\textbf{Operación lógica}: comprar una unidad con identidad 9F2A. \textbf{Filas}: pedido e inventario. \textbf{Tablas}: \texttt{purchases} e \texttt{inventory}. \textbf{Página}: bloque físico, por ejemplo P17. \textbf{Índice}: ruta secundaria, por ejemplo I5. \textbf{WAL}: registros que describen cambios y COMMIT. \textbf{Resultado idempotente}: estado que relaciona 9F2A con el pedido y la respuesta final.

Para que el retry sea seguro, el efecto de negocio y la marca de 9F2A deben confirmarse como una misma unidad transaccional. Si se descuenta stock y el proceso cae antes de persistir la clave idempotente, un retry puede volver a descontar. Si se registra la clave antes del efecto y luego se cae, el retry puede devolver éxito para una compra que no ocurrió. La atomicidad debe abarcar identidad, efecto y resultado.
\end{uaanswer}

\begin{uaexercise}
\textbf{4. Prediagnóstico razonado.}
\end{uaexercise}
\begin{uaanswer}
\begin{enumerate}
\item La compra confirmada debe reaparecer una sola vez. Más adelante se asocia con WAL, COMMIT durable, recovery e idempotencia.
\item El reporte debe usar una regla estable de visibilidad o declarar otra semántica. Se asocia con MVCC y snapshots.
\item Conviene evitar recorrer y ordenar filas irrelevantes. Se asocia con índice, planificador y \texttt{EXPLAIN}.
\item La cola crece a 30 MB/s mientras se mantenga el déficit. Se asocia con compactación, backlog y stall.
\item Un corte común elimina todas las copias en RAM. Replicación volátil no equivale a persistencia.
\end{enumerate}

La actividad es correcta porque primero formula comportamiento y falla; después introduce el nombre del mecanismo.
\end{uaanswer}

\begin{uaexercise}
\textbf{5. Timeline WAL y prefijo durable.}
\end{uaexercise}
\begin{uaanswer}
El orden recomendado es:

\[
\begin{aligned}
\text{request} &\rightarrow \text{preparar} \rightarrow \text{WAL UPDATE} \\
&\rightarrow \text{dirty pages} \rightarrow \text{COMMIT record} \\
&\rightarrow \text{flush hasta commitLSN} \rightarrow \text{ACK} \\
&\rightarrow \text{crash} \rightarrow \text{REDO}.
\end{aligned}
\]

El flush durable \textbf{no comienza} al agregar el primer \texttt{WAL UPDATE}. En ese paso comienza el append lógico al WAL y los bytes aún pueden estar en buffers. El flush ocurre después de agregar COMMIT y debe avanzar \texttt{flushLSN} hasta, como mínimo, \texttt{commitLSN}(T):

\[
\texttt{flushLSN} \geq \texttt{commitLSN}(T).
\]

Como WAL es secuencial, el prefijo incluye los registros de actualización anteriores que todavía no fueran durables. Si algunos ya lo eran, no hace falta reescribirlos; basta con avanzar la frontera persistida.
\end{uaanswer}

\begin{uaexercise}
\textbf{6. Dos reglas, dos fallas.}
\end{uaexercise}
\begin{uaanswer}
\textbf{Log antes que página}: el registro de REDO debe ser durable antes de permitir que la página modificada llegue al archivo de datos. Invertirlo puede dejar una página con un cambio que recovery no puede interpretar, completar o deshacer correctamente.

\textbf{COMMIT antes que ACK}: el registro de COMMIT debe ser durable antes de responder éxito. Invertirlo permite que el cliente observe una promesa que desaparece tras un crash.

La primera regla relaciona WAL y páginas; la segunda relaciona el motor y el cliente. Confundirlas conduce a la falsa idea de que una página debe estar persistida antes de cada ACK.
\end{uaanswer}

\begin{uaexercise}
\textbf{7. Matriz de crashes y reconstrucción completa.}
\end{uaexercise}
\begin{uaanswer}
\begin{center}
\small
\begin{tabular}{p{0.25\textwidth}p{0.20\textwidth}p{0.23\textwidth}p{0.21\textwidth}}
\toprule
Punto & Cliente & Recovery & Estado lógico \\
\midrule
Antes de COMMIT & No recibió éxito & Descarta o deshace trabajo incompleto & No confirmada \\
Después de append COMMIT, antes de flush & No recibió éxito & Depende de qué bytes sobrevivieron; el resultado es incierto para el cliente & Puede haber confirmado o no; retry con 9F2A resuelve \\
Después de flush, antes de ACK & No recibió éxito o recibe timeout & Conserva y rehace la transacción & Confirmada; retry devuelve mismo resultado \\
Después de ACK, antes de P17 & Recibió éxito & REDO materializa páginas todavía antiguas & Debe aparecer una sola vez \\
\bottomrule
\end{tabular}
\end{center}

El segundo y tercer caso muestran por qué un timeout no describe el estado remoto. La idempotencia durable transforma la incertidumbre del cliente en una respuesta estable.

La reconstrucción no depende del registro \texttt{COMMIT} aislado. Recovery combina:

\[
\text{páginas persistentes anteriores}
+
\operatorname{REDO}(\text{registros WAL confirmados pendientes}).
\]

Los registros \texttt{WAL UPDATE} describen cómo insertar la compra, descontar inventario y actualizar el índice. El registro \texttt{COMMIT} indica que ese conjunto de cambios debe formar parte de la historia confirmada. En síntesis: los registros de actualización contienen el \emph{cómo}; el COMMIT contiene la \emph{decisión}. Si se perdiese todo el dispositivo, haría falta un backup base más WAL archivado o una réplica apropiada.
\end{uaanswer}

\begin{uaexercise}
\textbf{8. REDO repetible.}
\end{uaexercise}
\begin{uaanswer}
Si todos los registros corresponden a P17, recovery omite LSN 118 y 120 porque la página ya declara \texttt{pageLSN}=120. Aplica 121 y luego 123, actualizando el \texttt{pageLSN} después de cada modificación.

Si recovery cae después de aplicar 121:

\begin{itemize}
\item si la página actualizada no llegó al medio, el siguiente recovery parte otra vez de la página durable con 120 y reaplica 121;
\item si llegó con \texttt{pageLSN}=121, el siguiente recovery omite 121 y continúa con 123.
\end{itemize}

Esto es idempotencia física del procedimiento. No evita que dos requests distintos creen dos compras; para eso se usa 9F2A como identidad de negocio.
\end{uaanswer}

\begin{uaexercise}
\textbf{9. Group commit cuantitativo.}
\end{uaexercise}
\begin{uaanswer}
Con un \texttt{fsync} por transacción se solicitan 600 sincronizaciones por segundo. Si fueran estrictamente seriales y cada una costara 2,4 ms, la demanda sería 1.440 ms de dispositivo por segundo y el máximo teórico sería aproximadamente:

\[
\frac{1}{0{,}0024} \approx 416{,}7\ \text{transacciones/s}.
\]

Con grupos completos de 12, el límite ideal sería 50 sincronizaciones/s y el costo físico amortizado:

\[
\frac{2{,}4\ \text{ms}}{12}=0{,}2\ \text{ms/transacción}.
\]

Sin embargo, con 600 transacciones/s llega una transacción cada 1,67 ms. Un timeout de agrupación de 1 ms no alcanza por sí solo para reunir 12; la ocupación real depende de concurrencia, cola existente y commits que lleguen mientras otra sincronización está en curso. La conclusión correcta es medir tamaño de grupo y latencia, no asumir grupos completos.

La garantía se conserva si cada ACK espera el flush que incluye su \texttt{commitLSN}.
\end{uaanswer}

\begin{uaexercise}
\textbf{10. Checkpoint, backup y réplica.}
\end{uaexercise}
\begin{uaanswer}
Usando unidades decimales:

\begin{center}
\begin{tabular}{lrr}
\toprule
Intervalo & Segundos & WAL máximo posterior \\
\midrule
2 min & 120 & 21.600 MB = 21,6 GB \\
5 min & 300 & 54.000 MB = 54 GB \\
15 min & 900 & 162.000 MB = 162 GB \\
\bottomrule
\end{tabular}
\end{center}

Más WAL posterior suele aumentar análisis y replay, por lo que presiona RTO. Checkpoints más frecuentes agregan E/S y pueden crear picos.

El checkpoint acota recuperación local. No protege contra pérdida total o corrupción del dispositivo: eso requiere backup y/o réplica durable. La réplica mejora disponibilidad y tolerancia a fallas de nodo, pero no reemplaza un backup frente a borrado lógico, corrupción propagada o retención histórica.
\end{uaanswer}

\begin{uaexercise}
\textbf{11. Timeline MVCC.}
\end{uaexercise}
\begin{uaanswer}
El timeline correcto es:

\begin{enumerate}
\item $v_1=120$ existe y está confirmada.
\item $T_R$ comienza y fija $S_R$.
\item La primera lectura elige $v_1$.
\item $T_W$ crea $v_2=135$ sin eliminar $v_1$.
\item $T_W$ confirma.
\item La segunda lectura de $T_R$ sigue usando $S_R$ y observa $v_1$ si el nivel conserva snapshot.
\item Un lector iniciado después del COMMIT puede observar $v_2$.
\item Al terminar $T_R$, $v_1$ queda reclamable cuando ningún otro snapshot la necesita.
\end{enumerate}

La justificación no es ``MVCC lo hace'', sino que ambas lecturas de $T_R$ usan la misma regla de visibilidad y $v_2$ no pertenecía a la historia observable al fijar $S_R$.
\end{uaanswer}

\begin{uaexercise}
\textbf{12. Snapshot y vista coherente.}
\end{uaexercise}
\begin{uaanswer}
Un snapshot no copia todas las páginas. Conserva metadatos que permiten decidir qué transacciones eran visibles. Una regla simplificada para una versión $v$ creada por $T_c$ y reemplazada por $T_r$ es:

\begin{quote}
$v$ es visible si $T_c$ estaba confirmada y visible para el snapshot, y $T_r$ no estaba confirmada/visible para ese mismo snapshot.
\end{quote}

Ejemplo incoherente: un reporte lee el total de pedidos antes de una compra y el total de stock después de esa misma compra. Las filas obedecen momentos distintos y el reporte puede presentar una igualdad contable imposible.
\end{uaanswer}

\begin{uaexercise}
\textbf{13. Transacción larga y bloat.}
\end{uaexercise}
\begin{uaanswer}
En 90 minutos se producen hasta:

\[
30.000 \times 90 = 2.700.000
\]

actualizaciones. No todas implican una versión retenida distinta, pero la cifra dimensiona la presión potencial. El snapshot antiguo puede obligar a conservar versiones que de otro modo serían reclamables.

Señales esperadas: edad de transacción, filas muertas, tamaño de tabla e índices, duración de VACUUM, cache hit rate descendente, más buffers por consulta, WAL creciente y p99 alto.

Intervenciones razonables: reducir alcance del reporte, usar réplica de lectura, dividir el trabajo, establecer timeout de transacción, exportar un snapshot controlado, revisar autovacuum y evitar sesiones idle in transaction. Matar la consulta sin diagnóstico puede ser necesario en emergencia, pero no es una estrategia de diseño.
\end{uaanswer}

\begin{uaexercise}
\textbf{14. MVCC no implica serializabilidad.}
\end{uaexercise}
\begin{uaanswer}
Invariante: al menos una guardia permanece activa. Estado inicial $A=1,B=1$.

\begin{itemize}
\item $T_1$ usa un snapshot con $A=1,B=1$ y escribe $A=0$.
\item $T_2$ usa el mismo estado visible y escribe $B=0$.
\item No existe conflicto write--write sobre una misma fila.
\item Ambas confirman y el resultado $A=0,B=0$ viola el invariante.
\end{itemize}

Mitigaciones: aislamiento serializable con detección de dependencias; lock explícito sobre el conjunto o una fila que materialice el invariante; rediseñar el dato para que ambas transacciones escriban un mismo recurso de coordinación. Cada alternativa paga bloqueo, abortos o complejidad de modelo.
\end{uaanswer}

\begin{uaexercise}
\textbf{15. De la visibilidad a la localización: árbol B y planes.}
\end{uaexercise}
\begin{uaanswer}
MVCC y el índice no son sustitutos. El árbol B conduce desde claves de búsqueda hacia filas o páginas candidatas; después, las reglas MVCC determinan si cada versión candidata es visible para el snapshot. Un recorrido conceptual es:

\[
\begin{aligned}
(8421,\texttt{CONFIRMED})
&\rightarrow \text{raíz}
\rightarrow \text{hoja ordenada por fecha}\\
&\rightarrow \text{referencias candidatas}
\rightarrow \text{verificación MVCC}.
\end{aligned}
\]

Plan A se lee desde \texttt{Seq Scan}: visita la tabla, filtra y descarta 9.980 filas, toca 840 buffers, ordena los resultados y finalmente entrega 20. El trabajo inicial es mucho mayor que la salida útil.

Plan B usa un \texttt{Index Scan} cuya clave coincide con filtros y orden. Toca 5 buffers en el ejemplo, encuentra las primeras 20 entradas útiles y \texttt{LIMIT} detiene el recorrido sin un \texttt{Sort} separado. Algunas entradas podrían no ser visibles para el snapshot; en ese caso, el recorrido continúa hasta reunir 20 resultados visibles.

La evidencia no es solo el tiempo: cambian acceso, filas descartadas, buffers y operador de orden. Los números son ilustrativos; deben reproducirse con datos, estadísticas y hardware reales.
\end{uaanswer}

\begin{uaexercise}
\textbf{16. Diseñar y cuestionar un índice compuesto.}
\end{uaexercise}
\begin{uaanswer}
La sentencia es:

\begin{verbatim}
CREATE INDEX idx_purchase_student_status_date
ON purchases(student_id, status, confirmed_at DESC);
\end{verbatim}

El prefijo de igualdades ubica estudiante y estado; la tercera columna entrega fechas en el orden requerido. El alto \emph{fan-out} significa que una página interna puede contener muchas claves separadoras y referencias. Por ello, la altura suele crecer lentamente y el motor necesita pocas lecturas de páginas para llegar a la hoja correcta.

\begin{itemize}
\item Solo por \texttt{status}: el índice puede ser poco útil porque la primera columna no está restringida. Puede evaluarse otro índice, uno parcial o un scan secuencial según cardinalidad.
\item Rango de fecha global: el orden del índice no comienza por fecha; un índice cuyo prefijo sea \texttt{confirmed\_at} podría ser mejor.
\item Por estudiante sin estado: puede aprovechar el prefijo \texttt{student\_id}, aunque quizá recorra varios estados.
\item Si 95\% está confirmado: \texttt{status} tiene baja selectividad; puede evaluarse un índice parcial, cambiar el orden o aceptar un scan secuencial.
\end{itemize}

En cada caso se observarían el nodo de acceso, filas estimadas y reales, filas descartadas, buffers, ordenamientos y tiempo. La decisión se valida con planes y workload, no con una regla sintáctica aislada.
\end{uaanswer}

\begin{uaexercise}
\textbf{17. El índice también cuesta.}
\end{uaexercise}
\begin{uaanswer}
Procedimiento:

\begin{enumerate}
\item medir uso de cada índice, consultas afectadas, p95/p99 y buffers;
\item medir inserts/s, WAL/s, espacio, cache hit y tiempo de mantenimiento;
\item reproducir el workload en staging con y sin índices;
\item capturar \texttt{EXPLAIN ANALYZE} de consultas candidatas;
\item retirar primero uno de manera reversible o marcarlo inválido/no usado según el motor;
\item observar una ventana suficiente y mantener DDL de rollback.
\end{enumerate}

Eliminar un índice puede mejorar escrituras, reducir WAL y liberar caché, pero empeorar una consulta rara y crítica. La decisión debe incorporar frecuencia y criticidad, no solo conteo de usos.
\end{uaanswer}

\begin{uaexercise}
\textbf{18. Ruta completa de una clave en una LSM-tree.}
\end{uaexercise}
\begin{uaanswer}
Escritura:

\[
(k,v) \rightarrow WAL \rightarrow memtable \rightarrow memtable\ inmutable
\rightarrow flush \rightarrow SSTable\ L0.
\]

Lectura: revisar memtables, seleccionar SSTables por rango, consultar Bloom, omitir las que demuestran ausencia, consultar índice y caché/bloque para las candidatas, y elegir la versión visible más reciente.

Una actualización crea otra versión. Una eliminación crea un tombstone. Ambos deben sobrevivir hasta que la compactación haya visto todos los runs donde podría existir una versión anterior; eliminar temprano el tombstone puede resucitar datos.

``Archivo imposible'' es incorrecto. La expresión precisa es: \emph{SSTable que puede omitirse para esta búsqueda porque el filtro asegura que la clave no está allí}. Bloom puede producir falsos positivos, pero no falsos negativos.
\end{uaanswer}

\begin{uaexercise}
\textbf{19. Diagnóstico cuantitativo de compactación.}
\end{uaexercise}
\begin{uaanswer}
Déficit:

\[
90-60=30\ \text{MB/s}.
\]

Con unidades decimales, faltan 9 GB para slowdown:

\[
9.000/30=300\ \text{s}=5\ \text{min}.
\]

Faltan 18 GB para stop:

\[
18.000/30=600\ \text{s}=10\ \text{min}.
\]

Son tiempos ideales: la capacidad y el ingreso pueden variar.

Hipótesis y evidencia:

\begin{itemize}
\item disco saturado: utilización, latencia y cola del dispositivo;
\item write amplification alta: bytes físicos / bytes lógicos y política de niveles;
\item demasiados archivos L0: conteo, frecuencia de flush y stalls por trigger;
\item CPU insuficiente para compaction: uso por threads y tiempo de compaction;
\item hot partition: bytes/QPS por shard;
\item snapshots/tombstones retenidos: edad, versiones y espacio no reclamable.
\end{itemize}

El stall es protección cuando limita temporalmente la admisión y permite estabilizar backlog. Es evidencia de incapacidad cuando se vuelve frecuente, prolongado o consume el SLO de escritura.
\end{uaanswer}

\begin{uaexercise}
\textbf{20. Dos cargas, dos arquitecturas defendibles.}
\end{uaexercise}
\begin{uaanswer}
\textbf{CampusShop.}

\begin{itemize}
\item $P$: confirmar compras e inventario y consultar historial.
\item $F$: crash, timeout, retry, pérdida de nodo y consulta lenta.
\item $G$: compra exactamente una vez a nivel lógico, inventario correcto, RPO cero frente a crash local confirmado y p99 de consulta acotado.
\item $S$: base transaccional con WAL, MVCC, índice compuesto, idempotency key 9F2A y replicación durable según disponibilidad requerida.
\item $T$: costo de coordinación, WAL, índices, VACUUM y failover.
\item Evidencia: crash injection, reintentos, \texttt{EXPLAIN ANALYZE}, p99, tiempo de recovery y replica lag.
\end{itemize}

\textbf{Telemetría.}

\begin{itemize}
\item $P$: ingerir 80.000 eventos/s y consultar rangos por dispositivo/tiempo.
\item $F$: crash, picos, pérdida de nodo, backlog, disco lento y late data.
\item $G$: pérdida máxima declarada, throughput sostenible, RTO y rangos consultables.
\item $S$: WAL, particionamiento temporal, LSM, SSTables, Bloom, compactación y retención.
\item $T$: write/space amplification, frescura eventual, espacio temporal y stalls.
\item Evidencia: soak test, p99, backlog slope, L0 files, WA, RPO/RTO y costo por GB.
\end{itemize}

La condición de revisión debe ser explícita. Ejemplo: abandonar la hipótesis LSM si el p99 de lectura puntual y el costo de compactación incumplen el SLO aun después de dimensionar y ajustar razonablemente.
\end{uaanswer}

\section{Criterio de autocorrección}

Una respuesta excelente:

\begin{enumerate}
\item define el objeto exacto de la garantía;
\item ubica la frontera física o lógica;
\item mueve la falla por la ejecución;
\item distingue mecanismo local de distribución;
\item reconoce qué no resuelve;
\item cuantifica o mide el costo;
\item propone evidencia que podría demostrar que estaba equivocada.
\end{enumerate}

\end{document}
